\section{Armar un panorama completo de capacidades a partir de tres tipos de evaluación}

\videofigure{figures/long-horizon-implications.jpg}{METR, GDPval y DeepScholar-Bench exponen, respectivamente, distintos cuellos de botella: el horizonte temporal, el producto económico y la síntesis de investigación.}{01:02:10--01:04:08}

\videofigure{figures/benchmark-realwork-gap.jpg}{Los tres tipos de benchmark presentan diferencias sistemáticas frente al trabajo real: calificación automática, tareas de una sola vez, ausencia de contexto o restricciones de fuentes.}{01:04:08--01:05:42}

\begin{table}[H]
\centering
\small
\begin{tabular}{>{\raggedright\arraybackslash}p{2.6cm}>{\raggedright\arraybackslash}p{3.4cm}>{\raggedright\arraybackslash}p{3.8cm}>{\raggedright\arraybackslash}p{3.5cm}}
\toprule
\textbf{Evaluación} & \textbf{Pregunta central} & \textbf{Principal ventaja} & \textbf{Principal punto ciego} \\
\midrule
METR & ¿Qué tan largas son las tareas que puede completar de forma estable? & Calibración uniforme en tiempo humano, permite rastrear tendencias & Sesgado hacia tareas de software, débil en calidad y valor \\
GDPval & ¿El producto supera al de un experto? & Trabajo real multimodal, interpretación económica directa & La preferencia de expertos es costosa, el contexto se simplifica \\
DeepScholar & ¿Puede realizar una síntesis de investigación confiable? & Evalúa a la vez recuperación, síntesis y citación & Dominios limitados, la exhaustividad sigue siendo difícil de juzgar automáticamente \\
\bottomrule
\end{tabular}
\end{table}

\videofigure{figures/known-unknown.jpg}{Tenemos más certeza de que el modelo progresa rápido en tareas cortas, claras y verificables, pero seguimos sin confianza en entornos de mucho contexto, entornos abiertos y recuperación exhaustiva.}{01:05:42--01:07:18}

\videofigure{figures/evaluation-takeaways.jpg}{El resumen del curso enfatiza la confiabilidad, la recuperación de errores, las características de la tarea y las métricas complementarias, en lugar de perseguir una sola tabla de posiciones.}{01:07:18--01:09:12}

\begin{importantbox}{La evaluación debe convertirse en la interfaz de entrenamiento del agent}
Las etiquetas de fallo de grano fino pueden guiar directamente la mejora: los errores de planeación alimentan los datos del planner, los errores de herramientas alimentan la tool policy, las fuentes omitidas alimentan el retrieval curriculum y las inconsistencias de citación alimentan el entrenamiento del verifier. Un buen benchmark no solo produce un ranking, también debe localizar la señal para la siguiente ronda de auto-mejora.
\end{importantbox}

\subsection{Preguntas abiertas}

\begin{itemize}
    \item ¿Cómo medir la colaboración organizacional a lo largo de días y entre personas, en lugar de una sola tarea en sandbox?
    \item ¿Cómo incorporar el contexto implícito, los permisos, la comunicación y las restricciones de cumplimiento en una evaluación reproducible?
    \item ¿Cómo evaluar la recuperación del agent tras un fallo, y no solo el éxito o fracaso final?
    \item ¿Cómo detectar el reward hacking, la contaminación de benchmarks (benchmark contamination) y el sesgo del evaluador (evaluator bias)?
    \item ¿Cómo combinar la duración de la tarea, la calidad, el costo, el consumo de energía y el riesgo en una frontera de Pareto que permita tomar decisiones?
\end{itemize}

\subsection{Resumen de la sección}

Ningún benchmark único puede representar a un agent de propósito general. Un panorama completo requiere dimensiones complementarias: el horizonte de la tarea, la confiabilidad, la calidad del producto, el costo económico, la obtención de contexto, la cobertura de fuentes y la recuperación de errores.

